\section{Multi-Agent Collaboration and Coordination}
\label{sec:rw-coordination}

Another line of work focuses on how specialized or persona-based agents coordinate their outputs. This issue applies both to agents responsible for different stages of the RE workflow and to agents representing different stakeholder or persona perspectives.

In MARE, coordination is primarily artifact-mediated. Specialized agents responsible for elicitation, modeling, verification, and specification exchange intermediate requirements artifacts through a shared workspace. When verification identifies problems, earlier agents can revise their outputs and trigger subsequent processing stages~\cite{mare2024}.

iReDev similarly adopts artifact-mediated coordination, but uses an event-driven communication mechanism. Its shared artifact pool stores intermediate and final requirements artifacts, while specialized agents continuously monitor changes and trigger subsequent actions accordingly~\cite{jin2025iredev}. This mechanism supports coordination across requirements-development stages without requiring unrestricted conversational interaction among all agents.

VibeIntent also coordinates specialized agents through structured intermediate artifacts. Visionary, Agenda, Interviewer, EndUser, Distiller, Sprint Manager, and Analyst agents progressively transform an initial project idea into Agile-ready requirements. In particular, the Distiller extracts requirements from different perspectives, consolidates related items, and adjudicates conflicts before the resulting requirements are transformed into backlog items~\cite{vibeintent2026}.

Coordination among persona agents can follow a different pattern. Social Agents deliberately prevents direct communication among its persona agents in order to preserve independent perspectives. Their outputs are instead combined through a moderator using an aggregation mechanism~\cite{bhattacharyya2026social}. Thus, coordination occurs at the aggregation level rather than through inter-agent deliberation.

MAJ-EVAL introduces explicit deliberation among stakeholder-aligned persona agents. Each agent first performs an independent evaluation, after which agents belonging to the same stakeholder group participate in a moderated free-form debate. An aggregator then combines qualitative feedback and quantitative judgments across stakeholder groups~\cite{chen2026majeval}. The framework therefore demonstrates how independently constructed stakeholder perspectives can be refined through multi-agent discussion before aggregation.

QUARE provides an explicit negotiation mechanism within requirements engineering. Its quality-specialized agents independently generate requirements, after which semantically overlapping or conflicting requirements are identified. Agents then engage in multi-round proposal, critique, and synthesis to resolve cross-quality conflicts before integration into a unified requirements model~\cite{cheng2026quare}. Unlike persona-oriented systems, however, the interacting agents represent software quality dimensions rather than human stakeholder perspectives.

These studies demonstrate several coordination patterns: artifact-mediated workflow collaboration, event-driven orchestration, independent aggregation, deliberative debate, and explicit negotiation. However, these mechanisms are usually developed either for RE-task agents or for persona-based agents. The systematic integration of stakeholder modeling with downstream multi-agent requirements development and coordination remains limited.
